% 课题研究目标
\begin{titlepage}
\makeatletter
\setcounter{page}{6}
\newpage
\AddToShipoutPictureBG{
    \begin{tikzpicture}[overlay,remember picture]
        \draw [line width=1pt]
        ($ (current page.north west) + (1.95cm,-1.5cm) $)
        rectangle
        ($ (current page.south east) + (-1.95cm,1.5cm) $);
    \end{tikzpicture}
% \hspace{1.8cm}\raisebox{30pt}{\textit{若所填内容超出此页可另加页，页码顺延。}}
}
\makeatother
\section{课题研究目标、研究内容、拟解决的关键问题}

\noindent\textbf{1. 研究目标}

\par
\indent
考虑到问题规模与约束复杂度的持续增长，传统群体智能与进化优化方法在约束引导能力、协同决策效率以及参数调控稳定性等方面逐渐暴露出局限性。特别是在大规模约束优化与复杂调度问题中，算法性能往往高度依赖人工经验设定，且不同问题之间缺乏统一的调控机制。围绕上述背景，本文工作以群体智能优化与学习型调控方法为分析对象，聚焦多智能体协同演化过程中约束调控与参数自适应这一共性核心问题，旨在从方法机理层面系统分析其作用方式与适用边界，为复杂约束环境下群体智能方法的智能化演进提供理论依据与方法支撑。

\noindent\textbf{2. 研究内容}

\par
\indent
围绕大规模约束优化与复杂调度问题中群体智能方法的演化行为与过程调控机制，已有研究表明，约束处理策略、协同结构设计以及参数调控方式是影响群体智能算法在复杂优化与调度场景中性能表现的关键因素。然而，这些因素多以经验规则或启发式方式引入，其内在作用机理及相互关系尚缺乏系统分析。基于此，本文工作从约束调控、群体协同与学习型自适应三个相互关联的视角出发，对群体智能优化方法在复杂约束与分布式决策场景中的行为特征与演化规律开展系统研究。

\noindent\textbf{2.1 多层协同惩罚调控机制设计}

\par
\indent
在复杂约束优化与调度问题中，惩罚机制不仅影响可行性引导方式，也直接作用于多智能体协同演化过程中的决策更新与收敛行为。本文工作关注惩罚调控在多智能体协同演化过程中的功能定位，旨在分析不同协同层级下惩罚机制对演化效率与系统稳定性的影响。具体而言，本部分重点研究：

\par
\indent
（1）惩罚机制在群体演化过程中对可行性引导与目标优化的调控作用及其随演化阶段变化的合理形式；

\par
\indent
（2）不同协同层级在信息共享、决策引导与稳定性维持中的功能分工及多层协同结构的组织原则；

\par
\indent
（3）统一评估与反馈机制在协同调控与学习型策略中的支撑作用，为后续协同调控与学习型策略提供一致的决策依据。

\noindent\textbf{2.2 智能通信策略：门控层次结构与变量选择机制}

\par
\indent
针对大规模分布式优化问题中通信成本高、变量耦合复杂所导致的协同效率下降现象，本文工作提出智能通信策略，通过多层门控机制与变量重要性评估相结合的方式，实现通信频率与通信内容的双维度优化。该策略旨在从时间维度（何时通信）和空间维度（通信哪些变量）两方面提升协同决策效率。具体而言，本部分重点研究：

\par
\indent
（1）\textbf{三层门控机制}：通过最小通信间隔约束、自适应冲突阈值判断以及阶段感知的概率采样，构建分层通信控制策略，在保证协同效果的前提下显著降低通信开销；

\par
\indent
（2）\textbf{变量重要性评估与选择}：基于冲突强度的指数移动平均（EMA）累积机制，动态评估各共享变量对协商效果的贡献度，采用Top-k选择策略仅对关键变量进行协商，降低单次通信的数据传输量；

\par
\indent
（3）\textbf{双维度协同优化}：分析门控机制与变量筛选在不同问题规模与约束强度下的协同作用规律，探讨两者的耦合效应及适用边界。

\noindent\textbf{2.3 基于强化学习的参数自适应调整机制}

\par
\indent
针对传统群体智能算法参数依赖人工调整、泛化能力弱以及难以适应动态演化环境的局限性，本文工作探索强化学习辅助的参数自适应调整机制，旨在实现演化过程中关键参数的智能化调控。本部分重点关注惩罚系数$\alpha$的自适应调整策略，同时探讨多参数协同调控的可行性与有效性。具体而言，本部分重点研究：

\par
\indent
（1）\textbf{轻量级RL策略网络设计}：针对分布式演化优化场景中状态空间简单、动作空间连续的特点，设计单层全连接策略网络，在保证调控效果的同时显著降低训练开销与计算复杂度；

\par
\indent
（2）\textbf{状态表征与奖励机制}：基于全局冲突指数（Global Conflict Index）与冲突变化趋势（$\Delta$CI）构建状态空间；奖励信号采用\textbf{相对改进率}以增强敏感性：
\[
r_t=\frac{f(x_t)-f(x_{t+1})}{|f(x_t)|}\times 100
\]
并据此构造目标动作，引导策略网络学习有效的参数调整规律；

\par
\indent
（3）\textbf{深度RL模型对比分析}：系统对比深度Q网络（DQN）、近端策略优化（PPO）、Actor-Critic（A2C）等主流深度强化学习模型与轻量级策略网络在参数效率、训练时间、内存占用及优化性能等方面的差异，为不同应用场景提供方法选择依据。

\par
\indent
强化学习辅助决策流程如图\ref{fig:rl_process}所示。

\begin{figure}[H]
\centering
\resizebox{0.9\linewidth}{!}{%
\begin{tikzpicture}[scale=0.85]
    % 状态
    \draw[thick] (0,3) rectangle (2.5,4.5);
    \node[align=center] at (1.25,3.7) {状态\\(State)};
    
    % RL Agent
    \draw[thick] (4,2.5) rectangle (6.5,4.5);
    \node[align=center] at (5.25,3.5) {RL Agent\\(策略网络)};
    
    % 动作
    \draw[thick] (8.5,3) rectangle (11,4.5);
    \node[align=center] at (9.75,3.7) {动作\\(Action)};
    
    % 奖励
    \draw[thick] (4,0.5) rectangle (6.5,2);
    \node[align=center] at (5.25,1.25) {奖励\\(Reward)};
    
    % 箭头
    \draw[->, thick] (2.5,3.5) -- (4,3.5);
    \draw[->, thick] (6.5,3.5) -- (8.5,3.5);
    \draw[->, thick] (9.75,3) -- (9.75,1.5) -- (6.5,1.5);
    \draw[->, thick] (4,1) -- (1.25,1) -- (1.25,3);
\end{tikzpicture}%
}
\caption{强化学习辅助决策流程}
\label{fig:rl_process}
\end{figure}

\noindent\textbf{3. 拟解决的关键问题}

\par
\indent
大规模约束优化与调度类问题通常具有变量规模大、约束结构复杂及决策过程高度耦合等特征，在群体智能求解过程中集中表现为解空间规模庞大、演化效率受限以及过程调控困难等突出挑战。为有效应对上述问题，有必要从解空间结构、决策行为引导以及过程调控机制等层面对关键问题进行凝练与分析。

\noindent\textbf{3.1 动态惩罚系数的自适应调整}

\par
\indent
传统惩罚函数法采用固定惩罚系数或简单规则式调整策略，难以适应不同搜索阶段的演化特征与问题特性。本文工作关注的关键问题包括：

\par
\indent
（1）如何根据当前种群的适应度分布与约束违反程度动态调整惩罚系数，实现可行性引导与目标优化的动态平衡？

\par
\indent
（2）如何避免惩罚系数调整过程中出现的过早收敛到不可行域或可行域边界的现象？

\par
\indent
（3）如何设计具有理论保证与实际效果的惩罚系数更新规则，确保在不同问题类型与约束强度下的稳健性？

\par
\indent
\textbf{解决思路：}采用基于强化学习的自适应调整机制，以全局冲突指数与冲突变化趋势作为状态表征，通过轻量级策略网络学习惩罚系数的动态调整规律，实现演化过程中约束处理强度的智能化调控。

\noindent\textbf{3.2 多智能体协同通信的效率优化}

\par
\indent
分布式多智能体优化过程中，频繁的信息交互与全变量协商导致通信成本高昂，严重制约算法的实际应用效率。本文工作关注的关键问题包括：

\par
\indent
（1）如何设计分层通信控制机制，在保证协同效果的前提下动态调整通信频率，避免不必要的通信开销？

\par
\indent
（2）如何有效评估共享变量的重要性，实现协商内容的智能筛选，降低单次通信的数据传输量？

\par
\indent
（3）如何平衡通信频率控制与变量选择策略，实现时间维度与空间维度的协同优化？

\par
\indent
\textbf{解决思路：}提出智能通信策略，通过三层门控机制实现通信频率的自适应控制，结合基于冲突强度EMA累积的变量重要性评估与Top-k选择策略，实现通信成本的双维度优化。

\noindent\textbf{3.3 强化学习与群体智能的深度融合}

\par
\indent
强化学习辅助参数自适应需要解决状态表征、奖励设计、训练效率以及与演化过程的耦合等问题。本文工作关注的关键问题包括：

\par
\indent
（1）如何提取能够有效反映当前演化状态与问题特征的状态表征，为策略网络提供充分的决策信息？

\par
\indent
（2）如何设计合理的奖励函数，引导RL策略网络学习到能够持续改善优化性能的参数调整策略？

\par
\indent
（3）如何在训练开销与优化性能提升之间取得平衡，避免过度训练导致的计算负担与轻量级设计导致的表达能力不足？

\par
\indent
\textbf{解决思路：}针对分布式演化优化场景中状态空间简单、动作空间连续的特点，设计单层策略网络，采用相对改进率作为奖励信号；同时系统对比多种深度RL模型，为不同应用场景提供方法选择依据。

\end{titlepage}
